%
% kontur-allgemein.tex -- Halbkreiskontur mit mehreren Einbuchtungen
%
% (c) 2026 Damien Flury, OST Ostschweizer Fachhochschule
%
\documentclass[tikz]{standalone}
\usepackage{amsmath}
\usepackage{times}
\usepackage{txfonts}
\usetikzlibrary{arrows,math,decorations.markings}
\definecolor{darkblue}{rgb}{0,0,0.7}
\definecolor{darkred}{rgb}{0.7,0,0}
\begin{document}
\def\skala{1}
\begin{tikzpicture}[>=latex,thick,scale=\skala,
	midarrow/.style={
		postaction={decorate},
		decoration={markings,mark=at position #1 with {\arrow[line width=1.2pt]{latex}}}}]

% Parameter der Zeichnung
\def\R{4.6}
\def\xa{-2.8}
\def\xb{-1.0}
\def\xm{2.6}
\def\e{0.45}

% Achsen
\draw[->] (-\R-0.6,0) -- (\R+0.7,0) coordinate[label={$\operatorname{Re}z$}];
\draw[->] (0,-0.6) -- (0,\R+0.7) coordinate[label={right:$\operatorname{Im}z$}];

% grosser Halbkreis C_R, im positiven Sinn durchlaufen
\draw[color=darkblue,line width=1.4pt,midarrow=0.5,midarrow=0.85]
	(\R,0) arc[start angle=0, end angle=180, radius=\R];
\node[color=darkblue] at (140:\R+0.42) {$C_R$};

% Strecken auf der reellen Achse, die ausgesparten Intervalle bleiben frei
\draw[color=darkblue,line width=1.4pt,midarrow=0.55]
	(-\R,0) -- (\xa-\e,0);
\draw[color=darkblue,line width=1.4pt,midarrow=0.55]
	(\xa+\e,0) -- (\xb-\e,0);
\draw[color=darkblue,line width=1.4pt,midarrow=0.6]
	(\xb+\e,0) -- (\xm-\e,0);
\draw[color=darkblue,line width=1.4pt,midarrow=0.55]
	(\xm+\e,0) -- (\R,0);

% Einbuchtungen: kleine Halbkreise oberhalb der Polstellen,
% jeweils im negativen Sinn durchlaufen
\draw[color=darkblue,line width=1.4pt,midarrow=0.5]
	(\xa-\e,0) arc[start angle=180, end angle=0, radius=\e];
\node[color=darkblue,font=\small] at (\xa,\e+0.36) {$C_{\varepsilon,1}$};
\draw[color=darkblue,line width=1.4pt,midarrow=0.5]
	(\xb-\e,0) arc[start angle=180, end angle=0, radius=\e];
\node[color=darkblue,font=\small] at (\xb,\e+0.36) {$C_{\varepsilon,2}$};
\draw[color=darkblue,line width=1.4pt,midarrow=0.5]
	(\xm-\e,0) arc[start angle=180, end angle=0, radius=\e];
\node[color=darkblue,font=\small] at (\xm,\e+0.36) {$C_{\varepsilon,m}$};

% Polstellen erster Ordnung auf der reellen Achse
\fill[color=darkred] (\xa,0) circle[radius=0.06];
\node[color=darkred] at (\xa,-0.46) {$x_1$};
\fill[color=darkred] (\xb,0) circle[radius=0.06];
\node[color=darkred] at (\xb,-0.46) {$x_2$};
\fill[color=darkred] (\xm,0) circle[radius=0.06];
\node[color=darkred] at (\xm,-0.46) {$x_m$};
\node[color=darkred] at (0.8,-0.46) {$\cdots$};

% Polstellen in der oberen Halbebene, vom Weg umlaufen
\fill[color=darkred] (-2.0,2.4) circle[radius=0.06];
\node[color=darkred] at (-2.0,2.4) [above right=-1pt] {$z_1$};
\fill[color=darkred] (0.8,1.3) circle[radius=0.06];
\node[color=darkred] at (0.8,1.3) [above right=-1pt] {$z_2$};
\fill[color=darkred] (2.6,2.4) circle[radius=0.06];
\node[color=darkred] at (2.6,2.4) [above right=-1pt] {$z_n$};
\node[color=darkred] at (1.8,1.95) {$\cdots$};

% Beschriftung der Integrationsgrenzen
\draw (-\R,-0.08) -- (-\R,0.08);
\node at (-\R,-0.46) {$-R$};
\draw (\R,-0.08) -- (\R,0.08);
\node at (\R,-0.46) {$R$};

\end{tikzpicture}
\end{document}
